\section{\MakeUppercase {Conclusion}}
\label{sec:Conclusion}
\renewcommand{\baselinestretch}{2.0}\normalsize


Resilience is a serious organizing concept in management and information systems research, increasingly central to how organizations articulate their response to disruption. Regulated sectors have built a mature preparedness backbone---continuity planning, operational-resilience regimes, threat-led testing---and the field has produced substantial conceptual, regulatory, and methodological work. Yet researchers and practitioners still struggle to show that resilience is more than declared: adaptive performance under disruption, and improvement across successive events, often remains harder to evidence than compliance itself \autocite{rota2024,skopik2016problem,woods2015four}.\\
\\
This review identifies three interlocking deficits---conceptual plurality, limited empirical validation, and methodological fragmentation---and responds with the Enactment--Scope framework and the research agenda developed in Sections~\ref{sec:Enactment-Scope}--\ref{sec:RD}. The implications for practice are sharper than accumulating further controls. Resilience is better treated as a continuous journey of learning and exploration within and across organizational boundaries than as a validated achievement certifiable through maturity scores or audit clearance. What requires validation is that journey itself: whether sensing, reconfiguration, and coordination improve as threat and disruption conditions evolve. Management already knows how to instrument adaptation to market conditions \autocite{teece2007explicating,duchek2019organizational}; the analogous task for security postures and coordination capabilities under evolving threats remains comparatively under-developed.\\
\\
For IS research, three imperatives follow: \textbf{define} adaptive performance before templates substitute for it; \textbf{measure} it from digital traces organizations already generate; and \textbf{govern} it through escalation, review, and cooperative activation rather than post-incident reporting alone. By advancing construct clarity, trace-based indicators, simulation under deep uncertainty, and governance designs that connect measurement to decision rights, information systems research is well positioned to help organizations validate adaptive capacity in digitally mediated environments---moving resilience from organizing principle to enacted, measurable, and revisable socio-technical capability.